\section[Introduction]{Introduction}

\subsection[Organisation du cours]{Organisation du cours}

\begin{frame}{Organisation du cours}
	Le cours est organisé autour de:
	\begin{itemize}
		\item Cours magistraux: 20 heures (2h x 10 sessions)
		\item Travaux dirigés: 16 heures (2h x 8 sessions)
		\item Travaux pratiques: 20 heures (2h x 10 sessions)
	\end{itemize}
	\vspace{0.5cm}
	Les créneaux horaires suivants seront utilisés:
	\begin{itemize}
		\item Mardi 15h45 - 17h45: TD / TP
		\item Mardi 18h - 20h: Contrôles
		\item Jeudi 13h30 - 15h30: Cours / TP
	\end{itemize}
\end{frame}

\begin{frame}{Évaluations}
	\begin{tabular}{|l|l|l|l|l|}
		\hline
		\textbf{Item} & \textbf{\%} & \textbf{Durée} & \textbf{Format} & \textbf{Date} \\
		\hline
		CC1 & 33\% & 20 min & En cours, QCM & 19/03 \\
		\hline
		TD \& TP & 33\% & - & Présence, complétion TP, projet & - \\
		\hline
		CC3 & 34\% & 1h30 & Conditions d'examen & 19/05 \\
		\hline
		CC4 & - & 1h30 & Conditions d'examen & 11/06 \\
		\hline
	\end{tabular} \\
	\vspace{0.5cm}
	Note: Le contrôle CC3 sera réalisé sur le créneau du mardi de 18h à 20h.
\end{frame}

\begin{frame}{Plan}
	\begin{itemize}
		\item Bases de l'apprentissage [6h]
		\begin{itemize}
			\item Introduction
			\item Typologie des problèmes d'apprentissage
			\item Problématiques de l'apprentissage artificiel
			\item Modèles linéaires
			\item Validation d'algorithme d'apprentissage
		\end{itemize}
		\item Algorithmes d'apprentissage [13h30]
		\begin{itemize}
			\item Algorithme Bayésien naïf [1h30]
			\item Analyse en composantes principales [2h]
			\item Arbres de décision [2h]
			\item Perceptrons et perceptrons multicouches [4h]
			\item K moyennes [2h]
			\item K plus proches voisins [2h]
			%\item Machines à vecteurs de support [4h]%
		\end{itemize}
	\end{itemize}
\end{frame}

\subsection[Paradigmes]{Paradigmes}

\begin{frame}{Paradigmes et systèmes de pensée}
	Daniel Kahneman (prix Nobel d'économie 2002) a théorisé 2 modes de fonctionnement du cerveau humain:
	\begin{itemize}
		\item \textbf{La pensée lente}: consciente, logique, nécessitant un effort
		\item \textbf{La pensée rapide}: inconsciente, intuitive, sans effort
	\end{itemize}
	\vspace{0.5cm}
	\begin{columns}
		\begin{column}{0.5\textwidth}
			\begin{figure}
				\centering
				\includegraphics[height=0.5\linewidth]{Images/velo_centrifuge}
				\label{fig:velo-centrifuge}
			\end{figure}
			%\vspace{-0.5cm}
			\centering
			\small Approche logique
		\end{column}
		\begin{column}{0.5\textwidth}
			\begin{figure}
				\centering
				\includegraphics[height=0.5\linewidth]{Images/velo-t-es-capable}
				%\caption[source: http://naitreetgrandir.com]{approche intuitive}
				\label{fig:velo-t-es-capable}
			\end{figure}
			%\vspace{-0.5cm}
			\centering
			\small Approche intuitive
		\end{column}
	\end{columns}
	\vfill 
	\tiny \textcolor{gray}{Source: Thinking, fast and slow - D. Kahneman - 2012 - Penguin}
\end{frame}

\begin{frame}{Paradigmes et systèmes de pensée}
	On peut établir un lien avec les paradigmes de programmation:
	\begin{itemize}
		\item La pensée lente $\rightarrow$ La programmation \textbf{déductive}: basée sur des règles logiques et explicites, issues des connaissances humaines. Le programme applique les règles aux données qui lui sont fournies.
		\item La pensée rapide $\rightarrow$ La programmation \textbf{inductive}: basée sur des données, l'expérience, les règles sont implicites et sont découvertes par l'algorithme. Le programme applique les règles qu'il a découvert à de nouvelles données, il généralise.
	\end{itemize}
\end{frame}

%\begin{frame}{Paradigmes et systèmes de pensée}
%	Un autre paradigme existe néanmoins:
%	\begin{itemize}
	%		\item Le mimétisme $\rightarrow$ La programmation \textbf{transductive}: il n'y a pas de règle, implicite ou explicite, indépendante des données d'apprentissage. L'algorithme ne généralise pas, il copie.
	%	\end{itemize}
%	\begin{figure}
	%		\centering
	%		\includegraphics[height=0.3\linewidth]{Images/mimétisme}
	%		\label{fig:mimétisme}
	%	\end{figure}
%	%\vspace{-0.5cm}
%	\centering
%	\small Approche mimétique
%\end{frame}

\subsection[Définitions]{Définitions}

\begin{frame}{Définitions de l'apprentissage automatique}
	\begin{block}{Définition 1}
		Champ d'étude qui donne aux ordinateurs la capacité d'apprendre sans être explicitement programmés. \\
		\hfill
		\textit{--- Arthur Samuel (1959)}
	\end{block}
	\vfill
	\tiny \textcolor{gray}{Source: Some Studies in Machine Learning Using the Game of Checkers - A.L. Samuel - 1959 - IBM Journal}
\end{frame}

\begin{frame}{Définitions de l'apprentissage automatique}
	\begin{block}{Définition 2}
		Un programme apprend d'une \textbf{expérience $E$} pour une \textbf{tâche $T$} avec une \textbf{performance $P$}, si $P$ s'améliore avec $E$. \\
		\hfill 
		\textit{--- Tom Mitchell (1997)}
	\end{block}
	%\pause
	\begin{itemize}
		\item \textbf{Expérience ($E$)} = Les données.
		\item \textbf{Tâche ($T$)} = Prédire, classer, regrouper, agir, générer.
		\item \textbf{Performance ($P$)} = La fonction de perte mathématique.
	\end{itemize}
	\vfill 
	\tiny \textcolor{gray}{Source: Machine Learning - T.M. Mitchell - 1997 - McGraw-Hill}
\end{frame}

\begin{frame}{Tout est nombre !}
	\begin{figure}
		\centering
		\includegraphics[width=1.0\linewidth]{Images/tout_est_nombre}
		\label{fig:tout-est-nombre}
	\end{figure}
\end{frame}
